% Lösungen Einheit 3 von 3 – Pythagoras in Figuren und Körpern (zusammenbau v0.1)
\einheitenkopf{Einheit 3 von 3 · Pythagoras in Figuren und Körpern}

\erg{21}{a) $9$ cm (Radius)}
%% TODO Lösung der Erklärzeile 22g) fehlt in der Bank
\erg{22}{a) Katheten: die zwei halben Diagonalen; Hypotenuse: die Rautenseite \quad b) halbe Grundseite $9$ cm; $h^2 = 41^2 - 9^2 = 1\,600$, $h = 40$ cm \quad c) halbe Grundseite $11$ m; $h^2 = 61^2 - 11^2 = 3\,600$, $h = 60$ m \quad d) halbe Grundseite $28$ cm; $h^2 = 53^2 - 28^2 = 2\,025$, $h = 45$ cm \quad e) halbe Grundseite $24$ mm; $h^2 = 30^2 - 24^2 = 324$, $h = 18$ mm \quad f) halbe Grundseite $30$ cm; $h^2 = 50^2 - 30^2 = 1\,600$, $h = 40$ cm \quad g) \ldots \quad h) halbe Grundseite $12$ cm; Schenkel$^2 = 1\,369$, Schenkel $= 37$ cm \quad i) Überstand $(50 - 32) : 2 = 9$ cm; Schenkel$^2 = 9^2 + 40^2 = 1\,681$, Schenkel $= 41$ cm \quad j) $h^2 = 5{,}3^2 - 2{,}8^2 = 20{,}25$, $h = 4{,}5$ m \quad k) $BD^2 = 3{,}9^2 - 3{,}6^2 = 2{,}25$, $BD = 1{,}5$ cm \quad l) Diagonale$^2 = 67{,}24$, Diagonale $= 8{,}2$ m \quad m) x-Unterschied $3$, y-Unterschied $4$; $PQ^2 = 25$, $PQ = 5$ \quad n) x-Unterschied $9$, y-Unterschied $12$; $EF^2 = 225$, $EF = 15$ \quad o) halbe Grundkante $9$ cm; $h_s^2 = 1\,681$, $h_s = 41$ cm \quad p) $r = 33$ cm; $s^2 = 33^2 + 56^2 = 4\,225$, $s = 65$ cm \quad q) $h^2 = 82^2 - 18^2 = 6\,400$, $h = 80$ cm \quad r) Grundkante halbieren; halbe Grundkante und Körperhöhe sind die Katheten, die Seitenhöhe ist die Hypotenuse: beide quadrieren, addieren, Wurzel ziehen. \quad s) Kegelhöhe $= \sqrt{23{,}04} = 4{,}8$ m; gesamt $18 + 4{,}8 = 22{,}8$ m \quad t) Bodendiagonale$^2 = 13$; Raumdiagonale$^2 = 13 + 6^2 = 49$, Raumdiagonale $= 7$ dm \quad u) Radius $= 7{,}2 : 2 = 3{,}6$ m; Radius und Dachhöhe sind die Katheten, die Stange (Mantellinie) ist die Hypotenuse: $s = \sqrt{3{,}6^2 + h^2}$ mit der Dachhöhe $h$. \quad v) Kegelhöhe $= \sqrt{6{,}2^2 - 3{,}9^2} = \sqrt{23{,}23} \approx 4{,}82$ m; Gesamthöhe $\approx 19{,}5 + 4{,}82 \approx 24{,}3$ m \quad w) Durchmesser $= 5{,}6$ cm; Diagonale $= \sqrt{11^2 + 5{,}6^2} \approx 12{,}3$ cm; Länge $\approx 12{,}3 + 4 \approx 16{,}3$ cm \quad x) $R(-2 | 3)$, $S(4 | -1)$; x-Unterschied $6$, y-Unterschied $4$; $RS^2 = 52$, $RS \approx 7{,}2$}
\erg{23}{a) Seil$^2 = 3{,}3^2 + 5{,}6^2 = 42{,}25$, Seil $= 6{,}5$ m}
\erg{24}{a) Nils hat die ganze Grundseite genommen; im Teildreieck steht nur die halbe, $0{,}7$ m. Richtig: $h^2 = 2{,}5^2 - 0{,}7^2 = 5{,}76$, $h = 2{,}4$ m \quad b) Die Höhe ist die Symmetrieachse und trifft die Grundseite in der Mitte; das rechtwinklige Teildreieck hat darum nur die halbe Grundseite als Kathete. \quad c) Diagonale des Kofferbodens $= \sqrt{4\,250} \approx 65{,}2$ cm; das ist mehr als $62$ cm – ja, schräg passt er. \quad d) Katheten $6$ (waagerecht) und $4$ (senkrecht), Hypotenuse AB}

24d)

\begin{ksys}[xmin=-3,xmax=5,ymin=-2,ymax=4,klein] \punkt{-2}{-1}{A} \punkt{4}{3}{B} \punkt{4}{-1}{} \end{ksys}

